% GENERATED FILE. Edit canonical lessons, then run npm run generate:books.
% canonical-source-hash: fnv1a64:34b5d53696b76ce3
% canonical-lessons: SA-C253-kitah, SA-C253-luta, SA-C253-vrscikah, SA-C253-godha, SA-C253-nakrah

\chapter{Insect, Spider, Scorpion, Lizard, Crocodile}
\label{ch:sa-insect-spider-scorpion-lizard-crocodile}
\hlchaptermodality{253}

\begin{chapteropening}
\textbf{By the end of this chapter:} \emph{I can say an insect, a spider, a scorpion, a lizard and a crocodile.}

Complete the last lesson of chapter 253: I can say an insect, a spider, a scorpion, a lizard and a crocodile.
\end{chapteropening}

\section[\texorpdfstring{\sk{कीटः} (kīṭaḥ)}{कीटः (kīṭaḥ)}]{\sk{कीटः} (kīṭaḥ) — an insect}

\label{lesson:SA-C253-kitah}



\begin{quote}
Before the new one: say the Sanskrit for jasmine, then the Sanskrit for holy basil.
\end{quote}

\subsection*{You'll want to know: \sk{कीटः}}
\textbf{\sk{कीटः}} — \emph{kīṭaḥ} — \textquotedblleft{}an insect\textquotedblright{}.

\begin{grammarlens}[title={What you've built}]
One more word for this chapter.
\end{grammarlens}

\subsection*{Your turn}
\begin{itemize}
\raggedright
  \item \emph{Say it:} \emph{kīṭaḥ}
  \item \emph{Say it:} \emph{kīṭaḥ}, once more
  \item \emph{Say it:} say \emph{kīṭaḥ}
  \item \emph{Recall:} say the Sanskrit for jasmine, then the Sanskrit for holy basil, then say \emph{kīṭaḥ} again
\end{itemize}

\begin{tcolorbox}[breakable,colback=teal!4,colframe=teal!35!black,title={Before you move on}]
What does \sk{कीटः} mean? (\textquotedblleft{}An insect\textquotedblright{}.) Say it once more.
\end{tcolorbox}
\section[\texorpdfstring{\sk{लूता} (lūtā)}{लूता (lūtā)}]{\sk{लूता} (lūtā) — a spider}

\label{lesson:SA-C253-luta}



\begin{quote}
Before the new one: say the Sanskrit for holy basil, then the Sanskrit for an insect.
\end{quote}

\subsection*{You'll want to know: \sk{लूता}}
\textbf{\sk{लूता}} — \emph{lūtā} — \textquotedblleft{}a spider\textquotedblright{}.

\begin{grammarlens}[title={What you've built}]
One more word for this chapter.
\end{grammarlens}

\subsection*{Your turn}
\begin{itemize}
\raggedright
  \item \emph{Say it:} \emph{lūtā}
  \item \emph{Say it:} \emph{lūtā}, once more
  \item \emph{Say it:} say \emph{lūtā}
  \item \emph{Recall:} say the Sanskrit for holy basil, then the Sanskrit for an insect, then say \emph{lūtā} again
\end{itemize}

\begin{tcolorbox}[breakable,colback=teal!4,colframe=teal!35!black,title={Before you move on}]
What does \sk{लूता} mean? (\textquotedblleft{}A spider\textquotedblright{}.) Say it once more.
\end{tcolorbox}
\section[\texorpdfstring{\sk{वृश्चिकः} (vṛścikaḥ)}{वृश्चिकः (vṛścikaḥ)}]{\sk{वृश्चिकः} (vṛścikaḥ) — a scorpion}

\label{lesson:SA-C253-vrscikah}



\begin{quote}
Before the new one: say the Sanskrit for an insect, then the Sanskrit for a spider.
\end{quote}

\subsection*{You'll want to know: \sk{वृश्चिकः}}
\textbf{\sk{वृश्चिकः}} — \emph{vṛścikaḥ} — \textquotedblleft{}a scorpion\textquotedblright{}.

\begin{grammarlens}[title={What you've built}]
One more word for this chapter.
\end{grammarlens}

\subsection*{Your turn}
\begin{itemize}
\raggedright
  \item \emph{Say it:} \emph{vṛścikaḥ}
  \item \emph{Say it:} \emph{vṛścikaḥ}, once more
  \item \emph{Say it:} say \emph{vṛścikaḥ}
  \item \emph{Recall:} say the Sanskrit for an insect, then the Sanskrit for a spider, then say \emph{vṛścikaḥ} again
\end{itemize}

\begin{tcolorbox}[breakable,colback=teal!4,colframe=teal!35!black,title={Before you move on}]
What does \sk{वृश्चिकः} mean? (\textquotedblleft{}A scorpion\textquotedblright{}.) Say it once more.
\end{tcolorbox}
\section[\texorpdfstring{\sk{गोधा} (godhā)}{गोधा (godhā)}]{\sk{गोधा} (godhā) — a lizard}

\label{lesson:SA-C253-godha}



\begin{quote}
Before the new one: say the Sanskrit for a spider, then the Sanskrit for a scorpion.
\end{quote}

\subsection*{You'll want to know: \sk{गोधा}}
\textbf{\sk{गोधा}} — \emph{godhā} — \textquotedblleft{}a lizard\textquotedblright{}.

\begin{grammarlens}[title={What you've built}]
One more word for this chapter.
\end{grammarlens}

\subsection*{Your turn}
\begin{itemize}
\raggedright
  \item \emph{Say it:} \emph{godhā}
  \item \emph{Say it:} \emph{godhā}, once more
  \item \emph{Say it:} say \emph{godhā}
  \item \emph{Recall:} say the Sanskrit for a spider, then the Sanskrit for a scorpion, then say \emph{godhā} again
\end{itemize}

\begin{tcolorbox}[breakable,colback=teal!4,colframe=teal!35!black,title={Before you move on}]
What does \sk{गोधा} mean? (\textquotedblleft{}A lizard\textquotedblright{}.) Say it once more.
\end{tcolorbox}
\section[\texorpdfstring{\sk{नक्रः} (nakraḥ)}{नक्रः (nakraḥ)}]{\sk{नक्रः} (nakraḥ) — a crocodile}

\label{lesson:SA-C253-nakrah}



\begin{quote}
Before the new one: say the Sanskrit for a scorpion, then the Sanskrit for a lizard.
\end{quote}

\subsection*{You'll want to know: \sk{नक्रः}}
\textbf{\sk{नक्रः}} — \emph{nakraḥ} — \textquotedblleft{}a crocodile\textquotedblright{}.

\begin{grammarlens}[title={What you've built}]
One more word for this chapter.
\end{grammarlens}

\subsection*{Your turn}
\begin{itemize}
\raggedright
  \item \emph{Say it:} \emph{nakraḥ}
  \item \emph{Say it:} \emph{nakraḥ}, once more
  \item \emph{Say it:} say \emph{nakraḥ}
  \item \emph{Recall:} say the Sanskrit for a scorpion, then the Sanskrit for a lizard, then say \emph{nakraḥ} again
\end{itemize}

\begin{tcolorbox}[breakable,colback=teal!4,colframe=teal!35!black,title={Before you move on}]
What does \sk{नक्रः} mean? (\textquotedblleft{}A crocodile\textquotedblright{}.) Say it once more.
\end{tcolorbox}
